\documentclass[9pt,a4paper]{article}
\usepackage[T1]{fontenc}
\usepackage[utf8]{inputenc}
\usepackage{tgtermes,tgheros,microtype}
\usepackage[a4paper,margin=15mm,headheight=18pt]{geometry}
\usepackage{xcolor,booktabs,tabularx,array,enumitem}
\usepackage[most]{tcolorbox}
\usepackage{fancyhdr,lastpage,titlesec}
\usepackage[hidelinks]{hyperref}
\usepackage{amsmath}

\definecolor{Navy}{HTML}{173753}
\definecolor{Teal}{HTML}{008F8C}
\definecolor{PaleTeal}{HTML}{E5F4F3}
\definecolor{Orange}{HTML}{F08A24}
\definecolor{PaleOrange}{HTML}{FFF1E2}
\definecolor{Slate}{HTML}{5B7083}
\newcommand{\sans}{\fontfamily{qhv}\selectfont}
\newcolumntype{Y}{>{\raggedright\arraybackslash}X}

\hypersetup{
  pdftitle={HeiQ XMembrane -- Test-cell minimum area and ERV scale check},
  pdfauthor={HeiQ XMembrane project team},
  pdfsubject={Coupon sizing and commercial ERV reasonableness check}
}

\pagestyle{fancy}
\fancyhf{}
\fancyhead[L]{\sans\footnotesize\color{Slate}HEIQ MATERIALS AG · XMEMBRANE}
\fancyhead[R]{\sans\footnotesize\color{Slate}TECHNICAL NOTE · 1 SEPTEMBER 2026}
\fancyfoot[L]{\sans\scriptsize\color{Slate}Conditional model and measurement-design evidence}
\fancyfoot[R]{\sans\scriptsize\color{Slate}Page \thepage\ of \pageref*{LastPage}}
\renewcommand{\headrulewidth}{0.45pt}
\titleformat{\section}{\sans\bfseries\color{Navy}\fontsize{14}{17}\selectfont}{\thesection}{0.6em}{}
\titlespacing*{\section}{0pt}{3mm}{1.2mm}
\setlength{\parindent}{0pt}
\setlength{\parskip}{0.8mm}
\setlist[itemize]{leftmargin=5mm,itemsep=0.7mm,topsep=1mm}
\renewcommand{\arraystretch}{1.15}

\newtcolorbox{mainbox}{enhanced,breakable,colback=PaleTeal,colframe=Teal,boxrule=0.7pt,arc=2mm,left=4mm,right=4mm,top=3mm,bottom=3mm}
\newtcolorbox{warningbox}{enhanced,breakable,colback=PaleOrange,colframe=Orange,boxrule=0.7pt,arc=2mm,left=4mm,right=4mm,top=3mm,bottom=3mm}

\begin{document}

\begin{tcolorbox}[enhanced,colback=Navy,colframe=Navy,boxrule=0pt,arc=0mm,left=8mm,right=8mm,top=5mm,bottom=5mm]
{\sans\bfseries\small\color{white}HEIQ XMEMBRANE PROJECT TEAM}\par
\vspace{2mm}
{\sans\bfseries\fontsize{20}{23}\selectfont\color{white}Test-cell minimum area and ERV scale check}\par
\vspace{1.5mm}
{\sans\color{white}Short technical note · 1 September 2026}
\end{tcolorbox}

\begin{mainbox}
\textbf{Decision.} Retain the planned 50$\times$50-mm maximum active window
(25~cm\textsuperscript{2}). The 40$\times$40-mm insert
(16~cm\textsuperscript{2}) is the smallest practical aperture that covers the
complete 0.5--5~A/m\textsuperscript{2} low-current screen. A
30$\times$30-mm insert (9~cm\textsuperscript{2}) is sufficient in the ideal
signal calculation for the 1.24--3.37~A/m\textsuperscript{2} range of the
selected direct-active-ERV case.
\end{mainbox}

\section{Minimum measurable active area}

The previous design already defined a 25-cm\textsuperscript{2} maximum aperture
and inserts down to 4~cm\textsuperscript{2}, but it did not state one explicit
minimum area. The updated calculation holds the total flow at 6~L/min per side
and requires a differential RH signal five times the combined expanded
uncertainty of two 0.1-RH-percentage-point channels.

\begin{tabularx}{\textwidth}{@{}Yrrr@{}}
\toprule
\textbf{Case} & \textbf{$j$} & \textbf{Ideal minimum} & \textbf{Insert} \\
 & \textbf{A/m\textsuperscript{2}} & \textbf{cm\textsuperscript{2}} & \textbf{cm\textsuperscript{2}} \\
\midrule
Inverse-design anchor, lower planned point & 0.50 & 13.09 & 16 \\
Inverse-design anchor & 1.00 & 6.55 & 9 \\
Selected direct-ERV profile minimum & 1.242 & 5.27 & 9 \\
Selected direct-ERV profile maximum & 3.366 & 1.94 & 4 \\
Order-unity reference & 10 & 65.5 & Does not fit at 6 L/min \\
Low-drag reference, $n=0.3$ & 10 & 218 & Does not fit at 6 L/min \\
\bottomrule
\end{tabularx}

The 25-cm\textsuperscript{2} aperture can measure the weaker references by
reducing flow: the ideal limits are 2.29~L/min for the order-unity case and
0.69~L/min for the low-drag case. Signal acquisition and high-flow film
characterization must therefore be separate modes.

\begin{warningbox}
These are gross Faradaic signal floors. Passive back-diffusion, gas films,
leakage, storage and electrode water can only reduce the net measured signal.
The physical coupon should therefore remain approximately 80$\times$80~mm to
provide sealing and guard land, even when a smaller active insert is used.
\end{warningbox}

\section{Consistency with a full-scale ERV}

The coupon and full-scale model are hydraulically consistent at the local
channel level.

\begin{tabularx}{\textwidth}{@{}YrrY@{}}
\toprule
\textbf{Quantity} & \textbf{Coupon} & \textbf{Direct ERV} & \textbf{Assessment} \\
\midrule
Channel gap & 2 mm & 2 mm & Direct match \\
Velocity & 0.5--4 m/s & 2 m/s & Full-scale point covered \\
Reynolds number at 2 m/s & 472 & 488/529 & Close laminar similarity \\
Flow length & 50 mm & 1.95/2.08 m & Coupon is a local test \\
Pressure drop & 5.8 Pa & 218/226 Pa & Consistent after length scaling \\
\bottomrule
\end{tabularx}

The full-drying case represents 150,000~CFM (70.79~m\textsuperscript{3}/s) and
34,425~m\textsuperscript{2}. RenewAire publishes membrane arrays up to
70,400~CFM, giving 2.13 arrays by airflow. Naples Community Hospital specified
approximately 249~Pa maximum static pressure for its twelve-core enthalpy array,
close to the model's 218--226~Pa. Kao Data London One uses thirteen
plate-exchanger units at 13.5~m\textsuperscript{3}/s each; the project flow is
5.24 such units and 40\% of that documented total.

\begin{mainbox}
\textbf{Scale conclusion.} Airflow, channel velocity, Reynolds number and
pressure drop are credible for a very large modular building installation.
The active membrane area remains an upper-end packaging challenge, not an
obvious physical contradiction. It must be distributed across multiple cores,
and commercial readiness cannot be claimed until electrodes, supports,
manifolds, leakage, cleaning and manufactured pressure drop are included.
\end{mainbox}

\textbf{Sources and reproducibility.}
Commercial references:
\href{https://renewaire.com/wp-content/uploads/2025/05/Submittal-SASeries-Sept2021-1.pdf}{RenewAire SA Series specification};
\href{https://renewaire.com/casestudies/naples-community-hospital/}{Naples Community Hospital case study}; and
\href{https://www.flaktgroup.com/en/references/kao-data-center/}{FläktGroup Kao Data Centre case study}.

The calculation is reproduced by
\path{scripts/design_eo_vapor_test_cell.py}; the complete table is in
\path{outputs/eo_vapor_cell_sizing_v1/minimum_detectable_active_area.csv}.

\end{document}
